\documentclass[10pt,a4paper]{article}
\usepackage[T1]{fontenc}
\usepackage{lmodern,geometry,amsmath,amssymb,booktabs,array,xcolor,fancyhdr,pgfplots}
\usepackage[colorlinks=true,urlcolor=prior,linkcolor=prior]{hyperref}
\pgfplotsset{compat=1.18}
\geometry{left=19mm,right=19mm,top=21mm,bottom=20mm,headheight=14pt}
\definecolor{ink}{HTML}{20363D}
\definecolor{prior}{HTML}{147D83}
\definecolor{posterior}{HTML}{C24931}
\definecolor{chainone}{HTML}{147D83}
\definecolor{chaintwo}{HTML}{C24931}
\definecolor{chainthree}{HTML}{8263A8}
\definecolor{chainfour}{HTML}{8C751F}
\pagestyle{fancy}\fancyhf{}
\fancyhead[L]{\small\color{ink}DSC--AC / Priors and posteriors}
\fancyhead[R]{\small 27 September 2026}
\fancyfoot[L]{\scriptsize Saved-run comparisons: convergence unverified}
\fancyfoot[R]{\small\thepage}
\setlength{\parindent}{0pt}\setlength{\parskip}{5pt}
\setlength{\emergencystretch}{2em}
\newcommand{\Normal}{\mathcal N}
\newcommand{\IG}{\operatorname{IG}}
\newcommand{\IW}{\operatorname{IW}}
\newcommand{\diag}{\operatorname{diag}}
\newcommand{\vecl}{\operatorname{vecl}}
\newcommand{\E}{\operatorname E}
\begin{document}
{\color{ink}\LARGE\bfseries Parameter priors and posteriors\par}
{\large The joint weekly dynamic stochastic correlation model\par}
\medskip
\textbf{What this document shows.} The prior specifies plausible parameter values before the
estimation likelihood is applied. The posterior updates that distribution using the observed
returns. This report gives the actual prior definitions, derives the conditional posterior
updates, and compares the priors with saved parameter draws from this repository.

\textbf{The current evidence is preliminary.} The latest saved estimation contains
\textbf{@@DRAWCOUNT@@ retained draws from four chains} (@@COUNTS@@).
The run stopped at its 12-hour time limit; its convergence record says \texttt{not\_checked}
because the diagnostic time budget was exhausted. The plots are therefore visualizations
of saved draws intended to approximate the posterior, not a validated posterior estimate.
No new MCMC run was launched to produce this document.

\section*{1. Which quantities are parameters?}
For 14 weekly percentage log returns, with no VAR lags ($p=0$), the observation model is
\begin{align*}
y_t\mid B_t,h_t,r_t &\sim\Normal(B_t,H_t),\\
H_t&=D_tP_tD_t,\qquad D_t=\diag\{\exp(h_{1t}/2),\ldots,\exp(h_{14,t}/2)\},\\
r_t&=\vecl(\log P_t)\in\mathbb R^{91}.
\end{align*}
The logarithm here is a \emph{matrix logarithm}; $r_t$ is not a vector of ordinary
correlations. The Archakov--Hansen inverse map produces the positive-definite correlation
matrix $P_t$ with unit diagonal [1].

\begin{tabular*}{\linewidth}{@{\extracolsep{\fill}}p{31mm}p{85mm}r@{}}\toprule
Quantity & Meaning & Distinct entries\\\midrule
$V$ & Covariance of weekly changes in expected returns $B_t$ & 105\\
$\sigma_{h,j}^2$ & Variance of changes in log variance, one per series & 14\\
$\sigma_{r,k}^2$ & Variance of changes in a matrix-log coordinate & 91\\\midrule
Fixed parameters & $\theta=(V,\sigma_h^2,\sigma_r^2)$ & 210\\
Weekly latent states & $B_t$, $h_t$, and $r_t$ & 119 per week\\\bottomrule
\end{tabular*}

\textbf{Interpretation.} Larger evolution variances allow less smooth latent paths.
The volatility parameter $\sigma_{h,j}^2$ is not the return variance $\exp(h_{jt})$.
Likewise, $V$ concerns changes in the mean, not the contemporaneous return covariance $H_t$.

\section*{2. What data are used?}
The first 104 returns, @@CALSTART@@ to @@CALEND@@, calibrate initial-state priors and are
excluded from the likelihood. The selected estimation likelihood covers \textbf{@@T@@ weeks},
@@MODELSTART@@ to @@MODELEND@@. Evolution-prior scales use rolling windows through that
estimation cutoff, including likelihood observations. This is empirical Bayes: the
hyperparameters are fitted to data and subsequently held fixed. Their calibration
uncertainty is not propagated. Historical state estimates are full-sample smoothers.

\clearpage
\section*{3. The priors actually used}
Let $N_0=104$ be the complete-row count in the reserved calibration block, $S_{\rm ML}$ its
maximum-likelihood covariance, and $S_0=.95S_{\rm ML}+.05\diag(S_{\rm ML})$.
The initial centers are
\[
\bar B=\bar y_0,\qquad \bar V_B=S_0/N_0,\qquad
\bar h=\log\diag(S_0),\qquad \bar r=\vecl\{\log\operatorname{Corr}(S_0)\}.
\]
\textbf{Mean states and their innovation covariance.}
\begin{align*}
B_1&\sim\Normal(\bar B,4\bar V_B),\qquad B_1\perp V,\\
B_t\mid B_{t-1},V&\sim\Normal(B_{t-1},V),\quad t=2,\ldots,T,\\
V&\sim\IW(\Psi_0,\nu_0),\quad
\Psi_0=N_0\bar V_B(0.01)^2=10^{-4}S_0,\quad\nu_0=104.
\end{align*}
The convention is $p(V)\propto |V|^{-(\nu_0+m+1)/2}
\exp\{-\tfrac12\operatorname{tr}(\Psi_0V^{-1})\}$, with $m=14$.
Thus $\E[V]=\Psi_0/89$, while the joint mode is $\Psi_0/119$ [2].
The diagonal marginal is $V_{jj}\sim\IG(45.5,\Psi_{0,jj}/2)$; the off-diagonal entries
are jointly constrained by positive definiteness and may be negative.

\textbf{Volatility states and their innovation variances.} For each $j=1,\ldots,14$:
\begin{align*}
h_{j1}\mid\sigma_{h,j}^2
&\sim\Normal(\bar h_j,11\sigma_{h,j}^2),\\
h_{jt}\mid h_{j,t-1},\sigma_{h,j}^2
&\sim\Normal(h_{j,t-1},\sigma_{h,j}^2),\quad t=2,\ldots,T,\\
\sigma_{h,j}^2&\sim\IG(a,b_h),\quad a=10,\quad b_h=@@HSCALE@@.
\end{align*}
\textbf{Correlation-coordinate states and their innovation variances.} For each $k=1,\ldots,91$:
\begin{align*}
r_{k1}\mid\sigma_{r,k}^2
&\sim\Normal(\bar r_k,11\sigma_{r,k}^2),\\
r_{kt}\mid r_{k,t-1},\sigma_{r,k}^2
&\sim\Normal(r_{k,t-1},\sigma_{r,k}^2),\quad t=2,\ldots,T,\\
\sigma_{r,k}^2&\sim\IG(a,b_r),\quad a=10,\quad b_r=@@RSCALE@@.
\end{align*}
Each block follows the same order as $B_t$: first-state prior, law of motion, innovation-variance prior.
Unlike $B_1\perp V$, the priors for $h_{j1}$ and $r_{k1}$ depend on their innovation variances.
Innovations are independent across dates and coordinates, conditional on their variances.
The innovation-variance priors are independent within and across the two families.
Here the \emph{shape--scale} density and its moments are
\[
p(s)=\frac{b^a}{\Gamma(a)}s^{-(a+1)}e^{-b/s},\quad s>0;\qquad
\E[s]=\frac b{a-1},\quad\operatorname{mode}(s)=\frac b{a+1}.
\]
The calibrated means are $\E[\sigma_h^2]=@@HMEAN@@$ and
$\E[\sigma_r^2]=@@RMEAN@@$. Each prior has coefficient of variation
$1/\sqrt{a-2}=0.354$. A common prior does not force the coordinate-specific posterior values
to be equal. The scale $b$ is not the prior mean.

\textbf{How we obtain the priors for $h_0$ and $f_0\equiv r_0$.}
This report denotes the correlation-coordinate state by $r_t$; equivalently, $f_0=r_0$.
From the same reserved 104-return block, form $S_0$ as above and set
\[
D_S=\diag\{\sqrt{S_{0,11}},\ldots,\sqrt{S_{0,14,14}}\},\qquad
C_0=D_S^{-1}S_0D_S^{-1}.
\]
Then the coordinate-specific initial-state priors are
\begin{align*}
h_{j0}\mid\sigma_{h,j}^2&\sim\Normal(\bar h_j,10\sigma_{h,j}^2),
&\bar h_j&=\log S_{0,jj},\\
r_{k0}\mid\sigma_{r,k}^2&\sim\Normal(\bar r_k,10\sigma_{r,k}^2),
&\bar r_k&=[\vecl(\log C_0)]_k.
\end{align*}
To calculate $\log C_0$, diagonalize $C_0=Q\Lambda Q'$ and use
$Q\diag(\log\lambda_1,\ldots,\log\lambda_{14})Q'$; retain its strict lower triangle
in column-major order. Calibration fixes the centers $\bar h$ and $\bar r$;
the pre-sample states $h_0$ and $r_0$ remain random, conditionally independent given
their innovation variances.

For either coordinate $z$ with innovation variance $s$ and prior center $\mu$, integrating
out $z_0$ gives $\operatorname{Var}(z_1\mid s)=10s+s=11s$, explaining the first-state priors above.
The path covariance is $\operatorname{Cov}(z_t,z_u\mid s)=s\{10+\min(t,u)\}$.
The marginal prior for $z_t$ is Student-$t$ with 20 degrees of freedom, location $\mu$,
and scale $\sqrt{b(10+t)/a}$; its standard deviation is $\sqrt{(10+t)b/(a-1)}$.

\textbf{How the evolution scales were calibrated.} For each rolling window, the code
differences rolling log variances and matrix-log coordinates, takes each coordinate's
sample variance, then the median across coordinates. The selected window is 104 weeks.
\begin{center}\small
\begin{tabular}{rrr}\toprule
Window (weeks) & Log-variance innovation mean & Coordinate innovation mean\\\midrule
@@CALIBRATION@@
\bottomrule\end{tabular}\end{center}

\clearpage
\section*{4. From the prior to the posterior}
Let $x=(B_{1:T},h_{1:T},r_{1:T})$ and let $\hat\lambda$ denote the fixed empirical-Bayes
hyperparameters. Bayes' rule gives
\[
p(\theta,x\mid y,\hat\lambda)\ \propto\
\underbrace{\prod_{t=1}^T\Normal(y_t;B_t,D_tP_tD_t)}_{\text{observation likelihood}}
\underbrace{p(x\mid\theta,\hat\lambda)p(\theta\mid\hat\lambda)}_{\text{state and parameter priors}}.
\]
With missing values the likelihood uses the observed subvector and covariance submatrix.
The desired parameter posterior is $p(\theta\mid y,\hat\lambda)=
\int p(\theta,x\mid y,\hat\lambda)\,dx$. This integration has no simple closed-form answer.
MCMC approximates it by sampling states and parameters repeatedly.

\textbf{Conditional posterior for $V$.} Set $\Delta B_t=B_t-B_{t-1}$. Conditional on
the sampled mean path,
\[
\boxed{V\mid B_{1:T},y,\ldots\sim
\IW\left(\Psi_0+\sum_{t=2}^T\Delta B_t\Delta B_t',\;\nu_0+T-1\right).}
\]
For the selected $T=@@T@@$, the degrees of freedom are @@POSTDF@@.
There are $T-1$ increments because the $B_1$ prior is independent of $V$.
The marginal posterior of $V$ averages these inverse-Wishart conditionals over uncertain
mean paths; it is generally a mixture, not a single inverse-Wishart distribution.

\textbf{Conditional posterior for either innovation variance $s$.}
The integrated random-walk density contributes $T$ Gaussian terms: the initial state
with variance $11s$ and $T-1$ changes with variance $s$. Therefore
\[
\boxed{s\mid z_{1:T},y,\ldots\sim\IG\left(a+\frac T2,\;
b+\frac12\left[\frac{(z_1-\mu)^2}{11}+
\sum_{t=2}^T(z_t-z_{t-1})^2\right]\right).}
\]
Use $(z,\mu,b)=(h_j,\bar h_j,b_h)$ or $(r_k,\bar r_k,b_r)$.
The shape is @@POSTSHAPE@@ for this sample. The marginal posterior again integrates
over uncertain paths; fitting one inverse-gamma density to all saved draws is unnecessary.

\textbf{How this matches the sampler.} The current \texttt{draw\_variances} function uses
an independence Metropolis--Hastings step. It proposes from
$\IG(a+(T-1)/2,b+\tfrac12\sum_{t=2}^T\Delta z_t^2)$ and accepts with probability
\[
\min\left\{1,\left(\frac{s'}s\right)^{-1/2}
\exp\left[-\frac{(z_1-\mu)^2}{22}\left(\frac1{s'}-\frac1s\right)\right]\right\}.
\]
That correction includes the initial-state density and targets the boxed conditional.
The proposal distribution alone is not the full conditional posterior.

\textbf{Conditional posteriors for the latent paths.}
Given $h,r,V$, $B_{1:T}\mid y,h,r,V$ is multivariate Gaussian and is sampled with a
Kalman filter and backward simulation. The $h_j$ and $r_k$ path conditionals are
non-Gaussian because they enter the covariance nonlinearly; their Gaussian random-walk
priors are combined with the observation likelihood using elliptical slice updates [3].
Correlations and volatilities are obtained draw by draw as $P_t=\mathcal C(r_t)$ and
$\exp(h_{jt}/2)$, before summarizing. In general, transforming a posterior mean is not
the same as taking the posterior mean of a transformed quantity.

\clearpage
\section*{5. Prior versus saved-draw distributions}
\textcolor{prior}{\textbf{Teal: exact prior density.}}\quad
\textcolor{posterior}{\textbf{Red: preliminary saved-draw density.}}
Each horizontal axis divides the parameter by its own prior mean; 1 is the prior mean.
All @@DRAWCOUNT@@ retained draws are pooled, so longer chains receive slightly more weight.
The red line is a kernel density estimate on the log scale, transformed back to this
positive scale with its Jacobian. These are marginal distributions, not state paths.
\par\medskip
@@DENSITIES@@
\par\medskip
\textbf{How to read a shift.} A red curve to the left describes smaller evolution
variances in the saved draws, hence smoother states. A narrower red curve describes less
dispersion in these draws. Neither feature proves correct mixing or convergence.
The two selected pairs illustrate familiar FX and equity labels; the appendix includes
every fixed parameter, so the examples are not the full evidence.

\clearpage
\section*{6. What the current comparison can establish}
The saved-draw means for all 14 log-variance innovation variances are
@@HMINRATIO@@--@@HMAXRATIO@@ times their common prior mean. Across the 91 correlation-coordinate
innovation variances, the corresponding ratios are @@RMINRATIO@@--@@RMAXRATIO@@.
These describe the current chains; they are not established posterior findings.
\begin{center}\small
\begin{tabular}{lrrr}\toprule
Example & Prior mean & Saved-draw mean & Ratio\\\midrule
@@COMPARISONS@@
\bottomrule\end{tabular}\end{center}
\textbf{Chain identity matters.} The traces below preserve each chain and show only retained
iterations after the configured 20-iteration burn-in. The colors denote chains
\textcolor{chainone}{1}, \textcolor{chaintwo}{2}, \textcolor{chainthree}{3}, and
\textcolor{chainfour}{4}. Similar pooled densities can conceal different chain locations.
\par\smallskip
@@TRACES@@
\par\smallskip
\textbf{Should the prior and posterior match?} There is no requirement that they match,
or that every posterior must be narrower. Similarity can mean weak information about a
parameter or agreement with the prior. A shift may reflect learning, prior conflict, or
sampling problems; its interpretation requires diagnostics and sensitivity checks.

\textbf{Before treating draw percentiles as credible intervals:} the repository's target
is four independent chains, rank-normalized split/folded $\widehat R<1.01$, bulk and tail
ESS at least 400, and mean MCSE/posterior-SD at most 0.05 for reported quantities [4].
The saved diagnostic record contains no usable values for these metrics. The 20-iteration
burn-in setting is a retention rule, not evidence that warm-up was sufficient. Passing
diagnostics would still leave prior sensitivity and model adequacy to assess.

\clearpage
\section*{7. Initial-state priors and reproducibility}
The following centers come from the reserved first 104 weekly returns.
The $B_1$ entries are percentage points of weekly expected return; its prior SD is
$2\sqrt{\bar V_{B,jj}}$. The log-variance centers use percentage-return units.
\begin{center}\small
\begin{tabular}{lrrr}\toprule
Series & $\bar B_j$ & Prior SD of $B_{j1}$ & $\bar h_j$\\\midrule
@@INITIALROWS@@
\bottomrule\end{tabular}\end{center}
The $B_1$ prior is multivariate, with covariance $4\bar V_B$, not 14 independent normals.
For $h_1$ and $r_1$, the marginal priors follow the Student-$t$ expression on page 2.
For $P_{ij,t}$ there is an induced joint prior obtained by mapping the entire $r_t$ vector;
assigning independent priors to ordinary pairwise correlations would change this model.

\textbf{Fixed-parameter smoothing is a different output.} When
\texttt{estimate\_parameters=false}, the workflow uses saved parameter means or saved
parameter draws while updating latent states. This is not a new estimation of
$p(\theta\mid y)$. Holding means fixed excludes parameter uncertainty; reusing saved
draws includes their variation but does not reweight them into an updated parameter
posterior when new observations arrive.

\textbf{Saved evidence.} Parameter file, estimated at @@ESTIMATEDAT@@:
\par{\footnotesize\nolinkurl{@@PARAMETERFILE@@}}
\par Run: {\footnotesize\nolinkurl{@@RUNDIR@@/estimation}}.
\par Source SHA-256: {\scriptsize\nolinkurl{@@SOURCEHASH@@}}.
The companion numerical CSV contains prior moments and retained-draw means and quantiles
for all 210 parameters. Its manifest records file hashes, chain counts, and validation
of positive definiteness, draw counts, and agreement with the saved parameter means.

\textbf{Implementation sources.}\par
{\footnotesize\begin{tabular}{@{}ll@{}}
\nolinkurl{toolbox/dsc_calibrate_priors.m} & Defines the priors\\
\nolinkurl{toolbox/dsc_prepare_inference.m} & Splits calibration and likelihood dates\\
\nolinkurl{toolbox/dsc_sample.m} & Implements the posterior updates\\
\nolinkurl{toolbox/dsc_save_parameters.m} & Preserves draws and chain IDs\\
\end{tabular}}
\par The generated LaTeX is self-contained and has no external figure dependencies.

\textbf{Methodological references.}\par
{\footnotesize
[1] Archakov and Hansen (2021). \emph{A New Parametrization of Correlation Matrices}.
\href{https://arxiv.org/abs/2012.02395}{Paper and reconstruction method}.\par
[2] MathWorks. \emph{iwishrnd}: inverse-Wishart parameterization and expectation.
\href{https://www.mathworks.com/help/stats/iwishrnd.html}{Official documentation}.\par
[3] Murray, Adams and MacKay (2010). \emph{Elliptical Slice Sampling}.
\href{https://proceedings.mlr.press/v9/murray10a.html}{PMLR 9, 541--548}.\par
[4] Vehtari et al. (2021). \emph{Rank-Normalization, Folding, and Localization}.
\href{https://arxiv.org/abs/1903.08008}{Bayesian Analysis 16, 667--718}.
}
@@APPENDIX@@
\end{document}
